\chapter{Storage and Persistence Boundaries}
\label{ch:storage}

\cnans{Microsoft::Xna::Framework::Storage} maps XNA's device-selection API onto one host
filesystem root. It provides usable save-file operations, but the abstraction is not a security
sandbox: two path inputs can escape their intended roots at the pinned revision.

\section{StorageDevice contract}

\cnaclass{StorageDevice} reports filesystem \texttt{FreeSpace}, \texttt{TotalSpace}, and
\texttt{IsConnected}. Its \texttt{Begin*}/\texttt{End*} pairs have XNA's asynchronous shape but
complete synchronously. Selector overloads accept player, free-space, and directory-count
constraints; CNA ignores the constraints because it has no removable-device choice.
\texttt{DeviceChanged} is never raised.

\texttt{SetAppNameEXT(appName)} selects the per-application root and invalidates any cached
resolution. Call it before storage or GamerServices access: achievements and leaderboards also
live below this root. \texttt{GetStorageRootEXT()} exposes the resulting path.

Storage result ownership is explicit. \texttt{EndShowSelector} returns
\texttt{std::unique\_ptr<StorageDevice>}; \texttt{EndOpenContainer} returns
\texttt{std::unique\_ptr<StorageContainer>}.

\begin{lstlisting}
StorageDevice::SetAppNameEXT("MyGame");

auto select = StorageDevice::BeginShowSelector(nullptr, nullptr);
std::unique_ptr<StorageDevice> device =
    StorageDevice::EndShowSelector(select.get());

auto opening = device->BeginOpenContainer(
    "MyGame Save Data", nullptr, nullptr);
std::unique_ptr<StorageContainer> container =
    device->EndOpenContainer(opening.get());
\end{lstlisting}

\section{Storage root resolution}

The root is resolved lazily, then cached until \texttt{SetAppNameEXT} changes the app name:

\begin{enumerate}
  \item \texttt{SDL\_GetPrefPath(nullptr, app)};
  \item \texttt{XDG\_DATA\_HOME/<app>} when SDL fails;
  \item \texttt{HOME/.local/share/<app>};
  \item the current working directory when the environment supplies no home.
\end{enumerate}

The final fallback is process-location-dependent and should not be treated as a stable save
location. If the root does not yet exist, space queries return \texttt{LLONG\_MAX}. For
\texttt{IsConnected}, CNA walks upward to the first existing ancestor; \texttt{true} therefore
does not prove that the per-app directory has been created.

\section{StorageContainer contract}

\cnaclass{StorageContainer} provides create, delete, existence, enumeration, and glob-pattern
operations for files and directories. Its \texttt{OpenFile} overloads return
\cnaclass{System::IO::Stream}. Paths are intended to be relative to the container.

\begin{lstlisting}
{
    std::unique_ptr<System::IO::Stream> out =
        container->CreateFile("savegame.dat");
    std::vector<SharpRuntime::bytecs> bytes = SerializeGameState();
    out->Write(bytes.data(), 0, static_cast<int>(bytes.size()));
}

if (container->FileExists("savegame.dat")) {
    std::unique_ptr<System::IO::Stream> in = container->OpenFile(
        "savegame.dat", System::IO::FileMode::Open);
    std::vector<SharpRuntime::bytecs> bytes(in->getLengthProperty());
    int offset = 0;
    while (offset < static_cast<int>(bytes.size())) {
        int n = in->Read(bytes.data(), offset,
                         static_cast<int>(bytes.size()) - offset);
        if (n == 0) break;
        offset += n;
    }
    if (offset != static_cast<int>(bytes.size())) {
        /* reject a truncated save */
    }
}
\end{lstlisting}

\texttt{Dispose()} is idempotent, sets \texttt{IsDisposed}, and raises \texttt{Disposing}; it
does not guard later operations. The container retains a non-owning
\texttt{const StorageDevice*}, so the device must outlive it. Treat a disposed container as
unusable even though the implementation still permits path operations.

\section{Open path-containment defects}

\begin{limitationnote}{Container name can escape the storage root}
\texttt{BeginOpenContainer(displayName)} retains the name verbatim. Construction joins
\texttt{root/displayName/playerFolder} and rejects only an empty name. An absolute name discards
the root under \texttt{std::filesystem::path::operator/}; \texttt{..} can climb above it.
\texttt{DeleteContainer} uses \texttt{ResolveContainedPath} and rejects these cases, so opening
and deleting do not share one boundary.
\end{limitationnote}

\begin{limitationnote}{File and directory names can escape a container}
The shared \texttt{ResolvePath(relative)} helper only evaluates
\texttt{storagePath/relative}. It rejects an empty string but neither absolute paths nor
\texttt{..}. For example, \texttt{CreateFile("../../outside.txt")} can write above the
container. Validate every externally influenced filename before passing it to this namespace.
\end{limitationnote}

The focused storage tests cover containment for \texttt{DeleteContainer}; they do not cover
\cnaclass{StorageContainer}, opening a container, enumeration, or the hand-written
\texttt{*}/\texttt{?} glob matcher. The untested boundary is therefore the same boundary where
the defects remain.

\section{File sharing is not enforced}

The fullest \texttt{OpenFile} overload accepts \cnaclass{System::IO::FileShare}, but comments
the parameter name out and constructs \cnaclass{FileStream} with path, mode, and access only.
The pinned \cnaclass{FileStream} has no share-mode constructor, and its
\texttt{std::fstream} implementation supplies no cross-process locking policy. Consequently,
\texttt{None}, \texttt{Read}, and \texttt{ReadWrite} produce the same CNA behavior. Do not use
\texttt{FileShare::None} as protection against concurrent save writers.

\section{Exception boundary and evidence}

\cnaclass{StorageDeviceNotConnectedException} is thrown when space queries translate a
filesystem exception. Container and file operations do not first check \texttt{IsConnected}
or uniformly translate host errors into this type.

The public headers and implementations establish ownership, root selection, ignored parameters,
and the two traversal paths. Tests establish the protected \texttt{DeleteContainer} cases and a
normal deletion. Container read/write behavior is source-proven here; the pinned storage test
set does not supply an end-to-end save round-trip or hostile container-path oracle. The example
uses fixed names and therefore does not demonstrate containment of untrusted input.
